% ============================================================
% preamble.tex — Style du dossier d'architecture
% Moteur : XeLaTeX (fontspec + polyglossia), TeX Live complet
% Polices : TeX Gyre Pagella (texte), TeX Gyre Heros (titres, tableaux),
%           DejaVu Sans Mono (code)
% ============================================================

% ---------- Métadonnées du document (utilisées partout) ----------
\newcommand{\docref}{DOC-ARCH-000}
\newcommand{\docver}{0.1}
\newcommand{\docstatus}{Brouillon}
\newcommand{\docdate}{2026-10-30}
\newcommand{\doctitle}{Dossier d'architecture — Cloud privé OpenStack}

% ---------- Langue et polices ----------
\usepackage{fontspec}
\usepackage{polyglossia}
\setmainlanguage{french}

\IfFontExistsTF{texgyrepagella-regular.otf}{%
  \setmainfont{texgyrepagella}[Extension=.otf,UprightFont=*-regular,BoldFont=*-bold,
    ItalicFont=*-italic,BoldItalicFont=*-bolditalic]%
}{\IfFontExistsTF{TeX Gyre Pagella}{\setmainfont{TeX Gyre Pagella}}{\setmainfont{Latin Modern Roman}}}

\IfFontExistsTF{texgyreheros-regular.otf}{%
  \setsansfont{texgyreheros}[Extension=.otf,UprightFont=*-regular,BoldFont=*-bold,
    ItalicFont=*-italic,BoldItalicFont=*-bolditalic,Scale=0.94]%
}{\IfFontExistsTF{TeX Gyre Heros}{\setsansfont{TeX Gyre Heros}[Scale=0.94]}{\setsansfont{Latin Modern Sans}}}

\IfFontExistsTF{DejaVuSansMono.ttf}{%
  \setmonofont{DejaVuSansMono}[Extension=.ttf,UprightFont=*,BoldFont=*-Bold,
    ItalicFont=*-Oblique,BoldItalicFont=*-BoldOblique,Scale=0.84]%
}{\IfFontExistsTF{DejaVu Sans Mono}{\setmonofont{DejaVu Sans Mono}[Scale=0.84]}{\setmonofont{Latin Modern Mono}[Scale=0.95]}}

\usepackage{unicode-math}
\IfFontExistsTF{texgyrepagella-math.otf}{\setmathfont{texgyrepagella-math.otf}}{\setmathfont{Latin Modern Math}}

\usepackage{microtype}
\usepackage{textcase}
\widowpenalty=10000
\clubpenalty=10000
\displaywidowpenalty=10000
\emergencystretch=2em
\raggedbottom

% ---------- Page ----------
\usepackage[a4paper,left=3.3cm,right=2.4cm,top=2.7cm,bottom=2.9cm,
  headheight=14pt,headsep=0.9cm,footskip=1.1cm]{geometry}
\setlength{\parindent}{0pt}
\setlength{\parskip}{0.55em plus 0.15em minus 0.1em}
\linespread{1.08}

% ---------- Couleurs : encre + un seul accent ----------
\usepackage{xcolor}
\definecolor{ink}{HTML}{1A1A1A}
\definecolor{muted}{HTML}{5F666D}
\definecolor{faint}{HTML}{8D939A}
\definecolor{hair}{HTML}{C4C8CC}
\definecolor{paper}{HTML}{F6F6F4}
\definecolor{accent}{HTML}{1F3A5F}
\definecolor{alert}{HTML}{8A2B2B}
\definecolor{ochre}{HTML}{8A6414}
\definecolor{moss}{HTML}{2E5A4B}
\AtBeginDocument{\color{ink}}

% Petites capitales espacées, pour les étiquettes (police sans)
\newcommand{\caps}[1]{{\sffamily\bfseries\addfontfeature{LetterSpace=9}\MakeUppercase{#1}}}
\newcommand{\hd}[1]{{\sffamily\bfseries\small\addfontfeature{LetterSpace=2}#1}}

% ---------- Titres ----------
\usepackage{titlesec}
\renewcommand{\chaptername}{Chapitre}
\titleformat{\chapter}[display]
  {\sffamily\bfseries\fontsize{22}{26}\selectfont\color{ink}}
  {\color{accent}\fontsize{46}{46}\selectfont\mdseries\thechapter}
  {2pt}
  {}
  [\vspace{9pt}{\color{hair}\titlerule}]
\titleformat{name=\chapter,numberless}[display]
  {\sffamily\bfseries\fontsize{22}{26}\selectfont\color{ink}}
  {}
  {0pt}
  {\vspace{41pt}}
  [\vspace{9pt}{\color{hair}\titlerule}]
\titlespacing*{\chapter}{0pt}{-4pt}{26pt}

\titleformat{\section}
  {\sffamily\bfseries\fontsize{13.5}{16}\selectfont\color{ink}}
  {\llap{\color{accent}\mdseries\thesection\hspace{1.1em}}}{0pt}{}
\titlespacing*{\section}{0pt}{22pt plus 4pt minus 2pt}{7pt plus 1pt}

\titleformat{\subsection}
  {\sffamily\bfseries\fontsize{10.8}{13}\selectfont\color{ink}}
  {\llap{\color{accent}\mdseries\thesubsection\hspace{1.1em}}}{0pt}{}
\titlespacing*{\subsection}{0pt}{15pt plus 3pt minus 2pt}{4pt plus 1pt}

\titleformat{\subsubsection}
  {\sffamily\bfseries\small\color{muted}}{\llap{\thesubsubsection\hspace{1.1em}}}{0pt}{}
\titlespacing*{\subsubsection}{0pt}{11pt}{3pt}

\setcounter{secnumdepth}{2}
\setcounter{tocdepth}{2}

% ---------- En-têtes et pieds de page ----------
\usepackage{fancyhdr}
\renewcommand{\chaptermark}[1]{\markboth{#1}{}}
\newcommand{\hdrstyle}{\sffamily\scriptsize\color{muted}\addfontfeature{LetterSpace=4}}
\pagestyle{fancy}
\fancyhf{}
\fancyhead[L]{\hdrstyle\MakeUppercase{\docref}\ \ \textperiodcentered\ \ v\docver\ — \docstatus}
\fancyhead[R]{\hdrstyle\nouppercase{\leftmark}}
\fancyfoot[L]{\hdrstyle Interne}
\fancyfoot[R]{\hdrstyle\thepage}
\renewcommand{\headrulewidth}{0.3pt}
\renewcommand{\footrulewidth}{0pt}
\renewcommand{\headrule}{\hbox to\headwidth{\color{hair}\leaders\hrule height \headrulewidth\hfill}}
\fancypagestyle{plain}{%
  \fancyhf{}
  \fancyfoot[L]{\hdrstyle Interne}
  \fancyfoot[R]{\hdrstyle\thepage}
  \renewcommand{\headrulewidth}{0pt}
}

% ---------- Tableaux ----------
\usepackage{colortbl}
\usepackage{booktabs}
\usepackage{array}
\usepackage{longtable}
\usepackage{xltabular}
\usepackage{multirow}
\usepackage{tabularx}
\renewcommand{\tabularxcolumn}[1]{>{\raggedright\arraybackslash}p{#1}}
\usepackage{etoolbox}
\setlength{\heavyrulewidth}{0.9pt}
\setlength{\lightrulewidth}{0.4pt}
\setlength{\aboverulesep}{0pt}
\setlength{\belowrulesep}{0pt}
\setlength{\tabcolsep}{7pt}
\renewcommand{\arraystretch}{1.32}
\AtBeginEnvironment{xltabular}{\sffamily\small\setlength{\parskip}{0pt}}
\AtBeginEnvironment{longtable}{\sffamily\small\setlength{\parskip}{0pt}}
\AtBeginEnvironment{tabularx}{\sffamily\small}
% Espace avant/après un tableau
\BeforeBeginEnvironment{xltabular}{\vspace{4pt}}
\AfterEndEnvironment{xltabular}{\vspace{6pt}}
\BeforeBeginEnvironment{longtable}{\vspace{4pt}}

% Marqueur « à compléter » dans les cellules : discret, lisible
\newcommand{\todo}{\leavevmode{\sffamily\bfseries\scriptsize\addfontfeature{LetterSpace=8}\color{ochre}À COMPLÉTER}}

% ---------- Listes ----------
\usepackage{enumitem}
\setlist[itemize]{leftmargin=1.3em,topsep=2pt,itemsep=1.5pt,parsep=0pt,labelsep=0.6em,
  label={\color{accent}\textendash}}
\setlist[enumerate]{leftmargin=1.6em,topsep=2pt,itemsep=1.5pt,parsep=0pt,labelsep=0.6em,
  label={\color{accent}\sffamily\small\bfseries\arabic*.}}

% ---------- Légendes ----------
\usepackage{float}
\usepackage{caption}
\gappto\captionsfrench{\renewcommand{\figurename}{Figure}\renewcommand{\tablename}{Tableau}}
\captionsetup{font={small},labelfont={sf,bf,color=accent},labelsep=quad,
  justification=raggedright,singlelinecheck=false,skip=8pt}
\captionsetup[figure]{position=below}
\captionsetup[lstlisting]{position=above,skip=6pt}

% ---------- Code ----------
\usepackage{listings}
\lstdefinestyle{base}{
  backgroundcolor=\color{paper},
  basicstyle=\ttfamily\footnotesize,
  keywordstyle=\color{accent}\bfseries,
  commentstyle=\color{faint}\itshape,
  stringstyle=\color{ochre},
  columns=fixed, keepspaces=true, basewidth=0.5em, sensitive=false,
  keywordstyle=[2]\color{muted},
  breaklines=true, frame=l, framerule=1.2pt, rulecolor=\color{accent},
  framesep=8pt, xleftmargin=14pt, framexleftmargin=8pt,
  numbers=left, numberstyle=\tiny\color{faint}\sffamily, numbersep=8pt,
  showstringspaces=false, tabsize=2, aboveskip=4pt, belowskip=6pt,
}
\lstdefinelanguage{pgsql}{
  sensitive=false,
  morekeywords={CREATE,TABLE,PRIMARY,KEY,NOT,NULL,UNIQUE,DEFAULT,CHECK,REFERENCES},
  morekeywords=[2]{VARCHAR,UUID,TIMESTAMPTZ,INTEGER,JSONB,TEXT,BIGSERIAL},
  morecomment=[l]{--}, morestring=[b]',
}
\lstset{style=base,language=pgsql}

% ---------- Encadrés : filet latéral, sans fond coloré ----------
\usepackage[most]{tcolorbox}
\newcommand{\boxlabel}[1]{{\sffamily\bfseries\scriptsize\addfontfeature{LetterSpace=8}\MakeUppercase{#1}}}
\tcbset{
  calloutbase/.style={
    enhanced, breakable, frame hidden, interior hidden, sharp corners,
    boxrule=0pt, left=11pt, right=0pt, top=3pt, bottom=3pt,
    before skip=9pt, after skip=9pt,
    fonttitle=\sffamily\bfseries\small, attach title to upper={\par\vspace{3pt}},
  },
}
\newtcolorbox{infobox}[1][]{calloutbase,
  borderline west={1.6pt}{0pt}{accent}, coltitle=accent, title={#1}}
\newtcolorbox{warnbox}[1][]{calloutbase,
  borderline west={1.6pt}{0pt}{alert}, coltitle=alert, title={#1}}
\newtcolorbox{okbox}[1][]{calloutbase,
  borderline west={1.6pt}{0pt}{moss}, coltitle=moss, title={#1}}
% À faire : filet pointillé, pour qu'on voie que c'est du provisoire
\newtcolorbox{todobox}[1][]{calloutbase,
  borderline west={1.6pt}{0pt}{ochre!85!white,dash pattern=on 3pt off 2pt},
  coltitle=ochre, title={\boxlabel{À compléter}\hspace{0.7em}#1}}

% ---------- Figures (TikZ) : traits fins, angles droits, une seule teinte ----------
\usepackage{tikz}
\usetikzlibrary{positioning,arrows.meta,calc,fit,backgrounds}
\tikzset{
  >={Stealth[length=5pt,width=4pt]},
  base/.style={font=\sffamily\footnotesize, text=ink, line width=0.6pt, draw=ink,
    align=center, inner sep=5pt},
  blk/.style={base, rectangle, fill=white, minimum height=1.0cm, minimum width=3.2cm},
  key/.style={blk, fill=accent!8},
  term/.style={blk, line width=1.1pt},
  flow/.style={line width=0.6pt, draw=ink},
  dflow/.style={flow, dash pattern=on 3pt off 2pt},
  lab/.style={font=\sffamily\scriptsize, text=muted, align=center, inner sep=2pt},
}

% ---------- Liens ----------
\usepackage{hyperref}
\hypersetup{
  colorlinks=true, linkcolor=accent, urlcolor=accent, citecolor=accent,
  pdftitle={\doctitle},
  pdfauthor={Équipe Infrastructure},
  pdfsubject={HLD, LLD et roadmap},
  bookmarksnumbered=true,
}
\usepackage{bookmark}

\usepackage{graphicx}
\graphicspath{{figures/}}
\newcommand{\isotoday}{\the\year-\ifnum\month<10 0\fi\the\month-\ifnum\day<10 0\fi\the\day}
\usepackage{siunitx}
\sisetup{locale=FR, group-separator={\,}, group-minimum-digits=4}
\usepackage{csquotes}

% ---------- Tables des matières ----------
\usepackage{titletoc}
\titlecontents{chapter}[0pt]
  {\addvspace{9pt}\sffamily\bfseries\small}
  {\contentslabel{2.3em}}
  {\hspace*{-2.3em}}
  {\hfill\contentspage}
\titlecontents{section}[2.3em]
  {\sffamily\small}
  {\contentslabel{2.4em}}
  {\hspace*{-2.4em}}
  {\color{hair}\titlerule*[5pt]{.}\color{ink}\contentspage}
\titlecontents{subsection}[4.7em]
  {\sffamily\footnotesize\color{muted}}
  {\contentslabel{2.9em}}
  {\hspace*{-2.9em}}
  {\color{hair}\titlerule*[5pt]{.}\color{muted}\contentspage}
\titlecontents{figure}[0pt]
  {\sffamily\small}
  {\makebox[2.8em][l]{\thecontentslabel}}
  {}
  {\color{hair}\titlerule*[5pt]{.}\color{ink}\contentspage}
